% Part: proof-theory
% Chapter: sequent-calculus
% Section: translations

\documentclass[../../../include/open-logic-section]{subfiles}

\begin{document}

\olfileid{pt}{seq}{tr}

\olsection{الترجمة بين~\Log{G1c} و\Log{G3c}}

ذكرنا أن~\Log{G1c} و\Log{G3c} كليهما سليم وكامل للمنطق الكلاسيكي؛ أي
إن كلًا منهما يبرهن جميع المتتاليات الصادقة في كل~!!{structure}، ولا
يبرهن سواها. ولذلك فهما يبرهنان المتتاليات نفسها. ويمكننا الآن إثبات هذه
الحقيقة بوسائل برهانية نظرية خالصة، من دون الالتفاف عبر مبرهنتي السلامة
والاكتمال. ويقدم مثل هذا البرهان \emph{تحويلات} لـ!!{proof}s في
\Log{G1c} إلى~!!{proof}s في~\Log{G3c} وبالعكس.

\begin{prop}
إذا كان~$\Log{G1c}\Proves S$، فإن~$\Log{G3c}\Proves S$.
\end{prop}

\begin{proof}
  نبين أنه إذا كان~$\pi$~!!a{proof} في~\Log{G1c}، فهناك~!!a{proof}
  $\pi'$ للمتتالية~$S$ في~\Log{G3c}. ونفعل ذلك بالاستقراء على
  $n=\pheight{\pi}$.

  إذا كان~$n=0$، كان~$\pi$ بديهية. والبديهية~$\lfalse\Sequent\quad$
  بديهية أيضًا في~\Log{G3c} (خذ السياقين~$\Gamma$ و$\Delta$ فارغين).
  أما~\Log{G1c} فيسمح ببديهيات~$!A\Sequent !A$ لأي~$!A$، في حين أن
  بديهيات الهوية في~\Log{G3c} ذرية فقط؛ لكننا أثبتنا في
  \olref[ex]{prop:G3c-axioms} أن جميع المتتاليات~$!A\Sequent !A$
  !!{provable} في~\Log{G3c}.

  إذا كان~$n>0$، انتهى~$\pi$ بقاعدة~$R$. ويضمن فرض الاستقراء وجود
  !!{proof}s لمقدمات تلك القاعدة في~\Log{G3c}؛ أي إن للـ!!{proof}s
  الجزئية المباشرة ترجمات في~\Log{G3c}. فإذا كانت~$R$ قاعدة في
  \Log{G3c} أيضًا، طبقناها على~!!{proof}s المترجمة للمقدمات. وتُحاكى القواعد الأربع
  التي تختلف بين النظامين على النحو الآتي: قبل تطبيق قاعدة~\Log{G3c}،
  نستخدم الإضعاف المقبول~\olref[adm]{prop:weak-G3c-adm} لإضافة الصيغة
  الفعالة الناقصة في حالتي~\LeftR{\land} و\RightR{\lor}، أو لإضافة
  الوقوع المحتفظ به للصيغة الرئيسة في حالتي~\LeftR{\lforall}
  و\RightR{\lexists}. أما إذا كانت~$R$ قاعدة إضعاف في~\Log{G1c}، فنطبق
  النتيجة نفسها عن قبول الإضعاف؛ وإذا كانت قاعدة تقلص، فنطبق
  \olref[inv]{prop:G3c-cont-adm}.
\end{proof}

\begin{prop}
إذا كان~$\Log{G3c}\Proves S$، فإن~$\Log{G1c}\Proves S$.
\end{prop}

\begin{proof}
  نمضي مرة أخرى بالاستقراء على~!!{height}~!!a{proof}~$\pi$ للمتتالية
  $S$. ويمكن~!!{prove} كل بديهية من بديهيات~\Log{G3c} في~\Log{G1c}
  انطلاقًا من بديهية ثم بعدة تطبيقات لقاعدتي~\LeftR{\Weakening}
  و\RightR{\Weakening}:
  \[
  \Axiom$!C \fCenter !C$
  \RightLabel{\LeftR{\Weakening}}
  \Deduce$!C, \Gamma \fCenter !C$
  \RightLabel{\RightR{\Weakening}}
  \Deduce$!C, \Gamma \fCenter \Delta, !C$
  \DisplayProof
\qquad
  \Axiom$\lfalse \fCenter$
  \RightLabel{\LeftR{\Weakening}}
  \Deduce$\lfalse, \Gamma \fCenter$
  \RightLabel{\RightR{\Weakening}}
  \Deduce$\lfalse, \Gamma \fCenter \Delta$
  \DisplayProof
  \]
  إذا انتهى~$\pi$ بقاعدة~$R$، أعطى فرض الاستقراء~!!{proof}s للمقدمات
  في~\Log{G1c}، ويمكننا تطبيق القاعدة~$R$ نفسها على تلك~!!{proof}s.
  والاستثناءات هي
  قواعد~\LeftR{\land} و\RightR{\lor} و\LeftR{\lforall}
  و\RightR{\lexists}، لأن صيغها تختلف بين النظامين. وتُشتق نسختا
  \Log{G3c} الأوليان في~\Log{G1c} هكذا:
  \[
  \Axiom$!A, !B, \Gamma \fCenter \Delta$
  \RightLabel{\LeftR{\land}}
  \UnaryInf$!A \land !B, !B, \Gamma \fCenter \Delta$
  \RightLabel{\LeftR{\land}}
  \UnaryInf$!A \land !B, !A \land !B, \Gamma \fCenter \Delta$
  \RightLabel{\LeftR{\Contraction}}
  \UnaryInf$!A \land !B, \Gamma \fCenter \Delta$
    \DisplayProof
\qquad
  \Axiom$\Gamma \fCenter \Delta, !A, !B$
  \RightLabel{\RightR{\lor}}
  \UnaryInf$\Gamma \fCenter \Delta, !A \lor !B, !B$
  \RightLabel{\RightR{\lor}}
  \UnaryInf$\Gamma \fCenter \Delta, !A \lor !B, !A \lor !B$
  \RightLabel{\RightR{\Contraction}}
  \UnaryInf$\Gamma \fCenter \Delta, !A \lor !B$
    \DisplayProof
  \]
  وتُشتق نسختا القاعدتين الكميتين على النحو الآتي:
  \[
  \Axiom$!A(t), \lforall[x][!A(x)], \Gamma \fCenter \Delta$
  \RightLabel{\LeftR{\lforall}}
  \UnaryInf$\lforall[x][!A(x)], \lforall[x][!A(x)], \Gamma \fCenter \Delta$
  \RightLabel{\LeftR{\Contraction}}
  \UnaryInf$\lforall[x][!A(x)], \Gamma \fCenter \Delta$
  \DisplayProof
\qquad
  \Axiom$\Gamma \fCenter \Delta, \lexists[x][!A(x)], !A(t)$
  \RightLabel{\RightR{\lexists}}
  \UnaryInf$\Gamma \fCenter \Delta, \lexists[x][!A(x)], \lexists[x][!A(x)]$
  \RightLabel{\RightR{\Contraction}}
  \UnaryInf$\Gamma \fCenter \Delta, \lexists[x][!A(x)]$
  \DisplayProof
  \]
\end{proof}

\end{document}
